% !TEX root = ../DoAn.tex
\documentclass[../DoAn.tex]{subfiles}
\begin{document}

\begin{center}
\Large{\textbf{TÓM TẮT NỘI DUNG ĐỒ ÁN}}\\
\end{center}

Trong bối cảnh các hệ thống mạng ngày càng mở rộng, lưu lượng trao đổi dữ liệu trở nên đa dạng, phức tạp và khó theo dõi bằng các phương pháp thủ công. Nhu cầu về một công cụ hỗ trợ quan sát, thu thập, phân tích và trình bày lưu lượng mạng một cách rõ ràng vì vậy ngày càng trở nên cần thiết, đặc biệt trong môi trường học tập, thực hành an toàn thông tin và điều tra kỹ thuật. Bên cạnh việc xem từng gói tin, người dùng còn cần một cách tiếp cận thuận tiện hơn để tổng hợp dữ liệu, theo dõi các luồng truyền thông, quan sát mối liên hệ giữa các thiết bị và khai thác thông tin phục vụ đánh giá hệ thống ngay trong cùng một môi trường làm việc.

Xuất phát từ thực tế đó, đồ án tập trung xây dựng một phần mềm phân tích gói tin mạng, hướng tới việc kết hợp nhiều nhóm chức năng trong cùng một hệ thống thống nhất. Phần mềm hỗ trợ bắt gói tin theo thời gian thực, đọc và ghi các tệp PCAP/PCAPNG, lọc và truy vết lưu lượng, thống kê dữ liệu, hiển thị bảng điều khiển phân tích, trực quan hóa sơ đồ mạng và cung cấp thêm lớp hỗ trợ phân tích bằng trí tuệ nhân tạo. Cách tiếp cận này giúp rút ngắn khoảng cách giữa dữ liệu gói tin thô và các thông tin có ý nghĩa, từ đó hỗ trợ người dùng quan sát, khám phá và đánh giá lưu lượng mạng một cách thuận tiện hơn.

Hệ thống được thiết kế nhằm phục vụ quá trình học tập, thực hành và khảo sát lưu lượng mạng một cách trực quan, có tổ chức và dễ tiếp cận hơn so với việc chỉ quan sát từng gói tin riêng lẻ. Thông qua việc liên kết dữ liệu ở mức gói tin với các góc nhìn tổng hợp ở mức luồng và toàn mạng, hệ thống giúp người dùng nhanh chóng nắm bắt bức tranh chung của phiên phân tích, xác định các luồng liên quan, theo dõi quan hệ giữa các thực thể mạng và khai thác dữ liệu phục vụ đánh giá kỹ thuật. Hệ thống phù hợp với nhu cầu học tập, thực hành phân tích mạng và hỗ trợ khảo sát kỹ thuật trong các tình huống cần quan sát dữ liệu một cách trực quan, có tổ chức và dễ theo dõi. Từ nền tảng đó, hệ thống cũng có thể tiếp tục được mở rộng trong các hướng nghiên cứu và phát triển về sau.

\end{document}
